\section{Setting and Authority Model}
\label{sec:model}

\subsection{The Setting as a Cyber-Physical System}
\label{sec:setting}

\Cref{fig:scenario} shows the scenario, a forest fire in hilly forest near a village,
taken from the scenario and system model we use with the rescue service
\pending{sign-off pending}. Three couplings make it a cyber-physical system rather than an
information system with mobile clients:
\begin{itemize}
  \item \emph{Orders actuate.} A decision moves crews and vehicles. An order that later
  turns out to be residual may already have been carried out; the system can label it,
  not undo it.
  \item \emph{Terrain decides connectivity.} Cellular coverage is good along the road and
  poor in the valleys; the mesh between vehicles is limited by range and terrain. Which
  devices can reach which follows from where people and vehicles are.
  \item \emph{Time without command is exposure.} While no one can record orders, crews act
  on the last orders they have, and new hazards go without a recorded response.
\end{itemize}

Requirements R1 and R2 follow from mobile-IT and incident-reporting findings that
responders need locally available information and durable later
documentation~\cite{rsg-kommunalplan,pilemalm2014enabling,westman2021mobilisation}; R3
follows from decision-provenance requirements and command
guidance~\cite{naja2010provenance,rsg-rad201}; R4 from the human-factors evidence that the
people who act need to know how certain their information is~\cite{endsley2017autonomy}.
\begin{itemize}
  \item \textbf{R1 (Offline operation).} Every device keeps working from previously
  synchronized data without continuous connectivity.
  \item \textbf{R2 (Durable submission).} Field submissions survive device restarts and
  prolonged partitions.
  \item \textbf{R3 (Attributable command under succession).} Every decision is attributable
  to an authenticated person and to the grant under which that person made it. For every
  decision the system determines whether the grant was valid under the organization's
  succession rule when the decision was made, and it keeps every decision for which it
  was not.
  \item \textbf{R4 (Visible certainty).} Every person who acts on a decision can see whether
  it stands, may still change, or was made without valid authority.
\end{itemize}

\subsection{People, Grants and Decisions}
\label{sec:system-model}

\textit{People and roles.} A primary commander A, a deputy B, a sector leader S, and a
superior officer C, the duty chief, who appoints commanders and may take over. Crews carry
out orders and report; operators at the regional centre follow the incident; reviewers
establish afterwards who decided what. The organization has one commander at a time.
Shared or contested authority, where several people must agree, is out of scope
(\cref{sec:thesis}).

\textit{Grants and succession acts.} C issues A a signed role grant. A succession rule, such
as ``B may assume command if A has been silent for $T$'', is signed by C in advance. A
handover from A to B, a takeover by C, an activation by B under a rule, and a promotion by an
operator are signed acts. Every grant and act is itself a record in the journal.

\textit{Decisions.} A decision is signed by the person who makes it and carries the act
under which that person believes they hold command, a per-person sequence number, and the
GNSS time of the device it was made on. Field submissions---observations, status reports,
requests---need no authority and are never invalid.

\textit{Devices.} Vehicle computers, handhelds and the regional server. A device exercises
the authority of the person using it and holds none of its own, so a commander who changes
device takes their authority along. Every device keeps the records it has, passes them on
over any link that is up, and sends heartbeats that carry its user's latest sequence
number. Records are never changed, merged or deleted, so they form a grow-only set and
duplicates and reordering do no harm.

\textit{Channels.} Cellular data to the regional server, a mesh between vehicles at the
scene, Wi-Fi between a handheld and a nearby vehicle, and Rakel voice radio. Voice is
outside the system, but it is how succession is usually agreed and how C tries to reach a
silent commander; the system records the outcome.

\textit{Partition cases.} \emph{P-A}: only A's device is cut off. \emph{P-B}: every vehicle
loses its cellular link, while the mesh still works. \emph{P-C}: the mesh also splits, so A
and B are in different parts. Independently, A may be \emph{fit} and answering voice, fit
but unreachable by voice, or \emph{unfit}---incapacitated, with or without a connected
device. Nobody can tell the last two apart from the network alone.

\textit{Assumptions.} Messages may be delayed, duplicated and reordered, and are eventually
resent. A device that holds a heartbeat holds every record its sender held when it sent it
(heartbeats travel with the records). Clock error between devices is bounded by $\epsilon$
while GNSS is available. Devices are not Byzantine, credentials are not stolen, and a person
signs from one device at a time.

\subsection{Who Commanded When, and the Status of a Decision}
\label{sec:statuses}

Every device computes who held command when from the authority records it holds, with one
rule: start from C's grant, apply the succession acts in timestamp order, and drop any act
its signer was not entitled to make at that moment (a handover by someone who no longer
held command, an activation under a rule whose primary had already handed over). Devices
that hold the same records therefore reach the same verdict. A decision then has one of
four statuses:
\begin{itemize}
  \item \textbf{Authoritative:} made by the person who held command at that time, and
  settled, so no record that could change this can still arrive (\cref{sec:settlement}).
  \item \textbf{Provisional:} valid on the records this device holds, but not yet settled.
  \item \textbf{Residual:} made by someone who did not hold command at that time, because
  their grant had been superseded or was never valid. Kept, attributed and labelled for
  review; never merged and never dropped.
  \item \textbf{In doubt:} made within the clock-error bound $\epsilon$ of a change of
  command, so the order is open. A person decides.
\end{itemize}
A residual decision is a \emph{legitimacy gap}. A valid decision that the server or the
successor does not yet hold, because its device has not reconnected, is a
\emph{completeness gap}. A crew can act on a provisional order that later becomes
residual; we call this a \emph{flip}, and count it as the clearest cost a policy imposes on
the people who act (\cref{sec:readers}). \pending{The status names and the display of
them to crews are to be checked in the expert walkthrough.}

The organization decides who commands, when command should move, and what to do about
residual decisions; the system records and enforces what it decided, applies its rules the
same way on every device, and shows every reader who commands and how certain that is. It
never decides who commands.

\subsection{The Authority Trade-off}
\label{sec:tradeoff}

Let A's device lose contact with every other device. It can receive no record, including
one saying that command has moved. While A is cut off, a design can provide at most two of:
(a) A keeps recording decisions; (b) command moves to B without A's participation;
(c) at most one device accepts decisions at any moment. With (a) and (b), A and B both
record; with (b) and (c), A must stop before B starts, on a local timer, which bounds (a);
with (a) and (c), B must wait for A. This is the availability--consistency tension of
partitioned systems~\cite{davidson1985consistency,gilbert2002cap,brewer2012cap} applied to
an authority role, and it holds by construction; we use it as framing, not as a result.
Succession is therefore a choice of policy. The question is what each choice costs, which
is the subject of the rest of the paper.
